\documentclass[11pt,a4paper]{article}

\usepackage[utf8]{inputenc}
\usepackage[T1]{fontenc}
\usepackage{amsmath}
\usepackage{geometry}
\usepackage{array}
\usepackage{booktabs}
\usepackage{hyperref}
\usepackage{longtable}

\geometry{margin=2.5cm}
\hypersetup{
    colorlinks=true,
    linkcolor=black,
    urlcolor=blue,
    pdftitle={Embedded NLP Under Microcontroller Constraints},
    pdfauthor={TODO}
}

\title{Embedded French Short-Text Classification Under Microcontroller Constraints}
\author{
    TODO: Author Name \\
    TODO: Institution / Program
}
\date{TODO: Submission Date}

\begin{document}

\maketitle

\begin{abstract}
This report documents a compact natural language classification pipeline designed for deployment on a microcontroller-class target, explicitly a Teensy 4.1 according to the project artifacts. The task is to classify short French workplace messages into 11 semantic categories while remaining compatible with the memory and execution constraints of an embedded system. Rather than using transformer-based architectures, the project relies on deterministic text normalization, hashed sparse lexical features, and a small multilayer perceptron (MLP) whose weights can be exported as C++ headers for embedded inference. The strongest compact checkpoint currently present in the workspace reaches a balanced accuracy of 91.65\% with 395{,}675 parameters, corresponding to roughly 1.51 MiB of float32 parameters stored in flash. This makes the project a useful case study in task-specific embedded NLP: domain-focused feature engineering and careful architecture search can provide strong classification quality without exceeding the practical constraints of a microcontroller deployment.
\end{abstract}

\section{Project Context}

The repository implements an embedded NLP classifier for short French messages. The end goal is not offline analysis on a workstation, but deployment on a microcontroller through direct export of model weights to a C++ header. The target hardware is identified in the project comments and generated files as a Teensy 4.1.

The classification problem is domain-specific. The model maps short texts to the following 11 labels:

\begin{center}
\texttt{ACCOUNTING, BANKING, BUSINESS, CYBER, GOSSIP, HR\_COMPLAINT, HR\_HIRING, INFRA, LOVE, MISC, TECH}
\end{center}

The design philosophy is clear from the code history:

\begin{itemize}
    \item avoid large language models and heavyweight inference stacks,
    \item use deterministic preprocessing that can be reproduced in C++,
    \item keep the model small enough for flash storage on the microcontroller,
    \item optimize balanced recall rather than only aggregate accuracy,
    \item verify that Python-side behavior and embedded-side behavior remain aligned.
\end{itemize}

\section{Hardware Constraints}

The Teensy 4.1 is built around a 600 MHz ARM Cortex-M7 and provides approximately 8 MiB of flash and 1 MiB of RAM according to the official PJRC specifications.\footnote{\url{https://www.pjrc.com/store/teensy41.html}} These numbers are large compared with older microcontrollers, but they remain highly constrained relative to conventional desktop or server NLP deployments.

For this project, the relevant constraints are:

\begin{itemize}
    \item \textbf{Flash budget}: model weights and static tables must fit into program memory.
    \item \textbf{RAM budget}: feature buffers, hidden activations, and runtime state must stay well below available RAM.
    \item \textbf{Implementation simplicity}: the feature extraction pipeline must be expressible in embedded C++ without dynamic vocabulary lookup.
    \item \textbf{Determinism}: preprocessing and hashing must match Python training behavior closely enough to keep predictions stable after export.
    \item \textbf{Latency predictability}: the system should avoid expensive tokenization models or runtime dependencies that are difficult to control on-device.
\end{itemize}

These constraints explain several architectural choices in the repository:

\begin{itemize}
    \item hashed features instead of an explicit learned embedding table,
    \item small MLP classifier instead of recurrent or transformer architectures,
    \item direct export of dense weights to \texttt{PROGMEM},
    \item bounded text length (up to 25 words),
    \item handcrafted, domain-oriented feature families rather than end-to-end representation learning.
\end{itemize}

\section{Dataset and Task Definition}

The final cleaned dataset tracked in the repository, \texttt{data/DataSetTeensyv9\_ULTRA\_CLEAN.csv}, contains 24{,}842 labeled samples. Table~\ref{tab:dataset} shows the class distribution.

\begin{table}[h]
\centering
\begin{tabular}{lr|lr}
\toprule
Class & Samples & Class & Samples \\
\midrule
ACCOUNTING & 2295 & HR\_COMPLAINT & 1817 \\
BANKING & 2176 & HR\_HIRING & 1909 \\
BUSINESS & 2461 & INFRA & 2895 \\
CYBER & 2740 & LOVE & 1954 \\
GOSSIP & 2343 & MISC & 1725 \\
TECH & 2527 & & \\
\bottomrule
\end{tabular}
\caption{Class distribution for the cleaned dataset used in the mature training stage.}
\label{tab:dataset}
\end{table}

An important characteristic of the project is that the dataset was not treated as fixed. The git history shows repeated cycles of:

\begin{itemize}
    \item generating synthetic examples,
    \item testing confusion points,
    \item cleaning templated artifacts,
    \item retraining,
    \item re-exporting the model for embedded use.
\end{itemize}

In other words, performance improvements came not only from model tuning, but from dataset shaping aimed at difficult boundaries such as \texttt{TECH} vs \texttt{CYBER}, \texttt{INFRA} vs \texttt{ACCOUNTING}, and \texttt{MISC} vs more semantically specific classes.

\section{Methodology}

\subsection{Text Processing}

The mature pipeline uses a deterministic analyzer that is deliberately easy to reproduce on embedded hardware:

\begin{itemize}
    \item accent removal by Unicode normalization,
    \item lowercasing,
    \item punctuation replacement,
    \item truncation to the first 25 words,
    \item generation of several feature families:
    \begin{itemize}
        \item character \(n\)-grams with word padding,
        \item word unigrams,
        \item word bigrams,
        \item word trigrams,
        \item positional markers for first and last words.
    \end{itemize}
\end{itemize}

These symbolic features are hashed into a fixed 8192-dimensional representation. This fixed-size design is critical for embedded deployment because it removes the need to store a dynamic vocabulary.

\subsection{Classifier}

The classifier is an \texttt{MLPClassifier} from scikit-learn with two hidden layers. The project uses Optuna to search over:

\begin{itemize}
    \item feature weights for each token family,
    \item character \(n\)-gram bounds,
    \item hidden layer sizes,
    \item activation function,
    \item regularization strength,
    \item learning rate.
\end{itemize}

The optimization objective emphasizes:

\begin{itemize}
    \item mean recall across classes,
    \item minimum class recall,
    \item balanced accuracy,
    \item control of overfitting,
    \item moderate model complexity.
\end{itemize}

This is a sensible objective for embedded classification: a single weak class can be more problematic than a small drop in global accuracy.

\section{Performance Under Resource Constraints}

\subsection{Checkpoint Comparison}

Table~\ref{tab:checkpoints} summarizes the most relevant checkpoints visible in the workspace. The table intentionally distinguishes between \emph{committed} results and later \emph{local experimental} results that are present in the working tree but not yet part of the branch history.

\begin{table}[h]
\centering
\footnotesize
\setlength{\tabcolsep}{4pt}
\begin{tabular}{>{\raggedright\arraybackslash}p{2.1cm} >{\raggedright\arraybackslash}p{1.9cm} c r r c c}
\toprule
Checkpoint & Status & Hidden Layers & Parameters & Float32 Size & Mean / Bal. Acc. & Min Recall \\
\midrule
2026-02-05 17:29 & Local run & \(96, 64\) & 793{,}451 & 3.03 MiB & 91.57\% & 82.61\% \\
2026-02-06 00:07 & Local run & \(48, 40\) & 395{,}675 & 1.51 MiB & 91.65\% & 84.06\% \\
2026-02-06 04:38 & Committed & \(80, 112\) & 665{,}755 & 2.54 MiB & 90.90\% & 83.77\% \\
\bottomrule
\end{tabular}
\caption{Relevant model checkpoints and their flash footprint.}
\label{tab:checkpoints}
\end{table}

Several observations follow directly from Table~\ref{tab:checkpoints}:

\begin{itemize}
    \item The best compact checkpoint is not the largest one.
    \item The \(48, 40\) network is the most interesting performance/resource trade-off in the workspace.
    \item Reducing parameter count from 793{,}451 to 395{,}675 halves the flash footprint while slightly improving both balanced accuracy and minimum recall.
    \item The latest committed result is not the strongest Pareto point; it is larger and less accurate than the compact local experimental checkpoint.
\end{itemize}

\subsection{Why the Compact Model Matters}

The strongest compact checkpoint (\(48, 40\) hidden units) deserves special attention:

\begin{itemize}
    \item Parameters: 395{,}675
    \item Float32 weight storage: 1{,}582{,}700 bytes (\(\approx\) 1.51 MiB)
    \item Mean recall / balanced accuracy: 91.65\%
    \item Minimum class recall: 84.06\%
\end{itemize}

On an 8 MiB flash device, storing 1.51 MiB of parameters consumes roughly 18.9\% of flash. This is substantial, but very manageable for a single embedded intelligence module. Even the largest visible checkpoint remains below half of the flash budget.

\section{Memory Budget Analysis}

\subsection{Flash Consumption}

The project stores model weights in flash via exported C++ headers. Since all parameters are emitted as float32 constants, the flash cost is essentially:

\[
\text{flash bytes} \approx 4 \times \text{parameter count}
\]

For the compact \(48, 40\) checkpoint:

\[
395{,}675 \times 4 = 1{,}582{,}700 \text{ bytes} \approx 1.51 \text{ MiB}
\]

This number excludes the rest of the firmware, feature extraction code, and minor metadata overhead, but the dominant cost is still the weight matrix from the 8192-dimensional input layer.

\subsection{RAM Consumption}

If the implementation materializes a dense float input vector of size 8192, then the input buffer alone requires:

\[
8192 \times 4 = 32{,}768 \text{ bytes}
\]

Adding hidden and output activations for the compact checkpoint gives a dense activation budget of approximately:

\[
(8192 + 48 + 40 + 11)\times 4 = 33{,}164 \text{ bytes} \approx 32.39 \text{ KiB}
\]

This is small enough to be practical on a device with 1 MiB RAM. Therefore, for the current design, \textbf{flash storage is the primary deployment constraint}, not activation RAM.

\subsection{Projected INT8 Savings}

An uncommitted refactor already contains an INT8 exporter with symmetric quantization. If the compact \(48, 40\) checkpoint were quantized without unacceptable accuracy loss, the theoretical parameter storage would fall from 1.51 MiB to about 0.377 MiB, a 4\(\times\) reduction. This is an important avenue for future work, but it should be described as a projection until end-to-end accuracy and parity are validated on the embedded side.

\section{Per-Class Behavior}

For the compact \(48, 40\) checkpoint, per-class recall remains above 90\% for most categories, with \texttt{MISC} as the weakest class at 84.06\%. This is expected. \texttt{MISC} acts as a residual category and therefore absorbs semantically heterogeneous utterances that do not strongly belong to one of the more specialized domains.

The highest per-class recalls in the compact checkpoint are:

\begin{itemize}
    \item ACCOUNTING: 95.86\%
    \item BANKING: 94.02\%
    \item BUSINESS: 93.09\%
    \item GOSSIP: 92.96\%
    \item HR\_HIRING: 92.41\%
\end{itemize}

This is a strong result given the deployment target. It indicates that the model is not merely fitting easy keywords, but has learned stable distinctions within a constrained feature space.

\section{What Enabled Strong Embedded Performance}

The main reason this project performs well under microcontroller constraints is that it does \emph{not} attempt to solve a general NLP problem with a general model. Instead, it narrows the problem in several useful ways:

\begin{itemize}
    \item the domain is controlled and label semantics are known,
    \item the text length is short,
    \item the feature extractor is highly task-specific,
    \item synthetic data generation was used to strengthen weak semantic boundaries,
    \item architecture search was constrained by deployment reality rather than benchmark prestige.
\end{itemize}

This is precisely the kind of trade-off embedded AI often requires: strong task performance through problem shaping and efficient representations, not brute-force model scale.

\section{Limitations}

The project also has important limitations that should be acknowledged in a research-style write-up:

\begin{itemize}
    \item The dataset is heavily synthetic and partially templated, even after cleaning.
    \item The repository records classification quality much more clearly than end-device latency, power, or throughput.
    \item Embedded parity is treated seriously in the code, but the repo does not yet include a full benchmark suite measured directly on the Teensy board.
    \item Some of the strongest checkpoints are present only as local workspace artifacts rather than committed experimental milestones.
\end{itemize}

Therefore, the report should present this work primarily as a successful \textbf{embedded NLP prototype aimed at deployment}, not yet as a finished hardware benchmark paper.

\section{Future Work}

The most credible next steps are:

\begin{itemize}
    \item validate inference latency on the Teensy 4.1 for representative sentence lengths,
    \item measure RAM usage directly in firmware rather than only estimating it,
    \item validate the INT8 quantized exporter on-device,
    \item compare dense 8192-float buffering with a more sparse or streaming feature accumulation strategy,
    \item add a fixed embedded test suite to prove Python/C++ prediction parity over time,
    \item enrich the evaluation with real user messages, not only synthetic examples.
\end{itemize}

\section{Conclusion}

This project demonstrates that domain-specific French message classification can achieve strong multi-class performance under genuine microcontroller constraints. The most interesting checkpoint present in the workspace reaches 91.65\% balanced accuracy with only 395{,}675 parameters and a float32 flash footprint of roughly 1.51 MiB. The result is significant because it is achieved without transformers, without a learned vocabulary table, and with an export path that is compatible with bare-metal C++ deployment on a Teensy-class device.

The broader lesson is that embedded NLP performance depends less on model scale than on disciplined system design: deterministic preprocessing, task-specific features, careful hyperparameter search, and explicit awareness of flash and RAM limits.

\section*{Notes for Finalization}

This draft intentionally leaves a few presentation details open:

\begin{itemize}
    \item author name,
    \item institution or course name,
    \item submission date,
    \item whether local experimental checkpoints should be included or whether the report should only cite committed branch artifacts.
\end{itemize}

\end{document}
